\section{Contexto del problema}
La diversidad del mundo natural es enorme: se estima que existen más de 5,000 especies de mamíferos, 10,000 de aves y 30,000 de peces, pero aún mayor, hay más de 400,000 tipos de flores \autocite{kagglePetalsMetal}. Distinguir entre ellas llega a requerir conocimiento especializado y, aún para una persona con experiencia, muchas flores son visualmente muy parecidas entre sí. Identificarlas de forma manual es lento y no es escalable cuando se tienen miles de fotografías. Por eso, resulta algo valioso poder contar con modelos que puedan reconocer el tipo de flor de forma automática a partir de una imagen. 

Este proyecto parte de la competencia "Petals to the Metal", encontrado en Kaggle, cuyo reto es construir un modelo que identifique el tipo de flor en un conjunto de imágenes. La competencia se limita a 104 clases de flores, y consiste en un problema de clasificación multiclase de imágenes en el que el modelo debe tener la capacidad de identificar valores finos como la forma, estructura y color para distinguir entre categorías cercanas. 

Este problema tiene varias dificultades que lo hacen relevante desde un enfoque de redes neuronales. Primeramente, las clases son heterogéneas, ya que algunas de ellas solo tienen un subtipo de flor y otras agrupan muchas subvariedades. Además, las imágenes de las flores no fueron capturadas en el mismo entorno, dando más variedad a las mismas. Por último, el número de imágenes por cada clase tampoco es uniforme, por lo que, considerando la métrica de evaluación, es un punto fuerte a tomar en cuenta que el modelo funcione de manera efectiva tanto para clases con pocos ejemplos como también para aquellas con muchos. 
